% ID: SEL-D0029-OB-E01
% ESTADO: cotejada. Segunda pasada visual completa y borrador compilado revisados.
% Fuente: Selectividad/matematicas_aplicadas/suplente septiembre/Examen-Matemáticas Aplicadas a las Ciencias Sociales.pdf
% SHA-256: 4b8884e307df4f5ddaf079e0f10be3d3f74d4793874eb03b355d34f4856579cb
% Página: 2; opción B; ejercicio original 1.

\textbf{(2.5 puntos)} Una fábrica de palas de pádel produce dos modelos A y B con los que obtiene un beneficio por cada pala de 30 y 20 euros respectivamente. Para la elaboración de una pala del modelo A se necesitan 90 g de fibra de carbono y 100 g de goma EVA, mientras que para una pala del modelo B son necesarios 100 g de fibra de carbono y 50 g de goma EVA. La fábrica dispone diariamente de 7.5 kg de fibra de carbono y 6.5 kg de goma EVA y quiere producir como máximo 60 unidades diarias del modelo A. Calcule cuántas palas de cada modelo tiene que fabricar para que el beneficio sea máximo y determine su importe.

\textquestiondown{}Sería posible una producción diaria de 49 palas del modelo A y 32 palas del modelo B?
